\section{Conclusion}
\label{sec:conclusion}

We framed LLM-based GPU performance reasoning as a software-engineering problem: deciding a kernel's \bandwidthbound/\computebound Roofline regime from software artifacts, before any profiling run, for developers and coding agents that cannot reach the target accelerator.
Off-the-shelf LLMs already handle this well for the FP32 and FP64 regimes from source alone, but sit at chance for FP16 until cross-compiled SASS is supplied---after which FP16 classification becomes near-perfect for the stronger models.
The decisive FP16 arithmetic is compiler-synthesized and invisible in source, so the source-versus-compilation boundary, not model choice, governs the hardest case; crucially, cross-compilation exposes it without any target-hardware access.

The signal is also practical to deploy.
It is cheap and tiered---a budget model triages the FP32 regime for a fraction of a cent per query with no GPU---stable enough under repeated queries to act on, with the most accurate model also the most consistent, and it beats lightweight static baselines in the high-value regimes.
No single configuration wins everywhere, so the right deployment is a cost-tiered cascade that matches the model and evidence tier to the target precision, not an LLM-only oracle.

These models do not replace profilers, they do not predict runtime in general, and numeric \rai is reliable only in identifiable regimes; the cascade itself remains a routing hypothesis we have not yet validated end-to-end.
What we show is narrower and useful: execution-free Roofline-regime classification is a well-posed software-engineering task, and off-the-shelf prompting already solves it in identifiable, cheaply-detectable evidence tiers.
